\section{Rechtsanwendung im Internet}

\subsection{Grundlagen}

\begin{definition}
    \textbf{DSGVO:} Die Datenschutz-Grundverordnung der EU regelt den Umgang und die Verarbeitung von Personen-bezogenen Daten in und ausserhalb der EU.
\end{definition}

\begin{remark}
    Die DSGVO findet Anwendung, wenn 
    \begin{itemize}
        \item Dienstleistungen in der EU angeboten werden
        \item Das Verhalten von Personen in der EU beobachtet wird (sofern dieses Verhalten in der EU stattfindet)
    \end{itemize}
\end{remark}

\begin{theorem}
    Jeder \textbf{Staat bestimmt selbst}, wo das eigene Recht Anwendung findet.
\end{theorem}

\begin{definition}
    \textbf{Anknüpfungspunkt:} Sachverhalt welchen Rechtsfall mit einem Land/Behörde in Beziehung setzt, respektive eine relevanz legitimiert.
\end{definition}

\begin{remark}
    Die blosse Abrufbarkeit einer Website in einem Land genügt in der Regel nicht als Anknüpfungspunkt; entscheidend ist meist, ob ein Angebot auf ein Publikum im Land ausgerichtet ist oder dort Auswirkungen hat.
\end{remark}

\begin{definition}
    \textbf{Souveränität:} Amtshandlungen auf Schweizer Boden sind grundsätzlich den Schweizer Behörden vorbehalten. Ausländische Herausgabebefehle auf z.B. Rechenzentren in der Schweiz können mit dem Schutz der Schweizer Souveränität über Art. 271 StGB (verbotene Handlungen für einen fremden Staat) kollidieren. 
\end{definition}

\subsection{Ermittlung des anwendbaren Rechts}

\begin{theorem}
    Die Ermittlung des anwendbaren Rechts passiert in Schritten:
    \begin{itemize}
        \item \textbf{Örtlich zuständige Behörde ermitteln:} \\
        Dies können Gerichte in mehreren Staaten gleichzeitig sein. \\
        Die Schweiz anerkennt und vollstreckt unter gewissen Bedingungen auch von ausländischen Gerichten ergangene Urteile.
        \item \textbf{Behörde bestimmt das anwendbare Recht nach eigenem Recht:} \\
        Schweizer Zivilgerichte wenden etwaige internationale ankommen sowie das Bundesgesetz über das internationale Privatrecht (IPRG) an um das anzuwendende Recht zu bestimmen.\\
        Schweizer Verwaltungs- und Strafbehörden wenden das Territorialprinzip an. 
    \end{itemize}
\end{theorem}

\begin{remark}
    \textbf{Rechtliche Vorkehrungen:} 
    \begin{itemize}
        \item Rechtswahl in B2B-Verträgen
        \item Günstiger Standort des eigenen Unternehmens (EU/nicht-EU/EU-Land)
        \item Disclaimer
        \item Abstimmung des Angebots auf die wichtigsten Rechtsordnungen
        \item Abstimmung des Angebots auf die strengsten Rechtsordnungen
        \item Günstige Wahl des Server-Standorts
        \item Günstige Wahl des Providers
    \end{itemize}
\end{remark}

\subsection{Beispiele}

\begin{example}
    \textbf{Konsumentenverträge:} Unterliegen strengen Sonderregeln zum Schutz von Konsumenten. Kunden können bei grenzüberschreitenden Online-Geschäften oft an ihrem Wohnsitz klagen und sich auf das Recht ihres Aufenthaltsortes berufen (z. B. deutsches Konsumentenschutzrecht für Kunden aus Deutschland)
\end{example}

\begin{example}
    \textbf{Unerlaubte Handlung:} Auf Schadenersatz geklagt werden kann in der Schweiz i.d.R. am Ort de Beklagten oder am Hanldungs- oder Erfolgsort.
    Falls der Täter je erwischt wird. Bis dahin verfolgen die Behörden am Ort des Opfers die Straftat.
\end{example}

\begin{example}
    \textbf{Urheberrecht:} Immaterialgüterrechte unterstehen grundsätzlich dem Recht des Staates, für den der Schutz beansprucht wird (Art. 110 IPRG).
\end{example}

\begin{example}
    \textbf{Aufsichtsrecht:} Unterliegt dem Territorialitätsprinzip bzw. dem Auswirkungsprinzip (z. B. Verbot von Alkoholwerbung für Jugendliche in der Schweiz).
\end{example}
